\chapter{Abstract}
\label{ch:abstract}

Branch-and-bound algorithms are traditionally viewed as exact methods for solving NP-hard optimization problems, often limited in practice by their exponential worst-case complexity.
Recent theoretical work has shown that, under suitable conditions, standard branch-and-bound frameworks can exhibit Polynomial-Time Approximation Scheme (PTAS) behavior for parallel machine scheduling problems.

This project provides an empirical investigation of these results, focusing on the practical impact of profiling-based optimizations proposed for the uniform machines setting.
The study concentrates on the identical parallel machine scheduling problem, a special case of uniform machines, in order to evaluate these optimizations in a controlled and interpretable environment.

A branch-and-bound algorithm was implemented and extended with approximation-aware stopping criteria and profiling-based prioritization and pruning mechanisms.
A flexible experimental pipeline was developed to generate instances, execute large-scale experiments, and collect detailed performance data across different problem sizes and approximation parameters.

The experimental results show that profiling-based strategies can substantially reduce redundant exploration of the search space and enable high-quality approximate solutions to be obtained within practical computational limits.
Moreover, the analysis highlights the strong influence of instance structure on algorithmic behavior, suggesting that a careful characterization of instance hardness is essential for understanding the practical performance of approximation-oriented branch-and-bound methods.
